The main strategy of the current solution is correct; the bug report is also fair that some standard convex-geometric facts were used without being written out. I will keep the same proof idea, but make those points explicit.

Throughout, indices are taken modulo \(n\). We assume \(s_1,\dots,s_n\) are listed in \emph{counterclockwise} order; if they are listed clockwise, reverse all labels. For a nonzero vector \(u\) and a compact convex set \(K\subset \mathbb R^2\) I write
\[
F_K(u):=\{x\in K:\langle u,x\rangle=h_K(u)\}
\]
for the exposed face of \(K\) with outer normal \(u\).

I will use the standard planar convention:

- a **vertex** of a polygon is a \(0\)-dimensional face;
- an **edge** of a polygon is a nondegenerate \(1\)-dimensional face.

So once we know an exposed face is a singleton, it is a vertex; and once we know it is a nondegenerate segment, it is an edge.

Set
\[
E_i:=[s_i,s_{i+1}],\qquad a_i:=s_{i+1}-s_i,\qquad R:=P+Q,\qquad x_i:=s_i+t_i .
\]

---

## 1. Preliminary facts

### 1.1. \(P\) and \(R\) are genuine polygons

Because \(S\) is in convex position and \(n\ge 3\), the points of \(S\) cannot all be collinear: if they were, every point of \(S\) except the two extreme ones on that line would lie in the convex hull of the others, contradicting convex position. Hence \(P=\operatorname{conv}(S)\) is \(2\)-dimensional.

Also, every vertex of \(\operatorname{conv}(S)\) belongs to \(S\), and every point of \(S\) is extreme because \(S\) is in convex position. Therefore the vertices of \(P\) are exactly
\[
s_1,s_2,\dots,s_n
\]
in the given cyclic order, and the edges of \(P\) are exactly
\[
E_i=[s_i,s_{i+1}].
\]

Next, for any finite sets \(A,B\subset \mathbb R^2\),
\[
\operatorname{conv}(A)+\operatorname{conv}(B)=\operatorname{conv}(A+B),
\qquad A+B:=\{a+b:a\in A,\ b\in B\}.
\]

#### Proof
First, \(A+B\subset \operatorname{conv}(A)+\operatorname{conv}(B)\), and the right-hand side is convex, so
\[
\operatorname{conv}(A+B)\subseteq \operatorname{conv}(A)+\operatorname{conv}(B).
\]

Conversely, let
\[
x\in \operatorname{conv}(A)+\operatorname{conv}(B).
\]
Then
\[
x=\sum_i \alpha_i a_i+\sum_j \beta_j b_j
\]
for some \(a_i\in A\), \(b_j\in B\), with \(\alpha_i,\beta_j\ge 0\), \(\sum_i\alpha_i=\sum_j\beta_j=1\). Hence
\[
x=\sum_{i,j}\alpha_i\beta_j(a_i+b_j),
\]
and the coefficients \(\alpha_i\beta_j\) are nonnegative and sum to
\[
\sum_{i,j}\alpha_i\beta_j=\Bigl(\sum_i\alpha_i\Bigr)\Bigl(\sum_j\beta_j\Bigr)=1.
\]
So \(x\in \operatorname{conv}(A+B)\). Thus
\[
\operatorname{conv}(A)+\operatorname{conv}(B)=\operatorname{conv}(A+B).
\]
\(\square\)

Applying this with \(A=S\), \(B=T\), we get
\[
R=P+Q=\operatorname{conv}(S+T).
\]
Since \(S+T\) is finite, \(R\) is the convex hull of finitely many points.

Also \(T\neq \varnothing\) (both signs occur), so choose \(t_0\in T\). Then
\[
P+t_0\subseteq P+Q=R.
\]
Because \(P+t_0\) is \(2\)-dimensional, \(R\) is \(2\)-dimensional as well. Therefore \(R\) is a genuine convex polygon.

---

### 1.2. Exposed faces of a convex hull of finitely many points

Let \(A\subset \mathbb R^2\) be finite and \(K=\operatorname{conv}(A)\). Define
\[
M_A(u):=\{a\in A:\langle u,a\rangle=\max_{x\in A}\langle u,x\rangle\}.
\]
Then
\[
F_K(u)=\operatorname{conv}(M_A(u)).
\]

#### Proof
Let
\[
m:=\max_{x\in A}\langle u,x\rangle .
\]
If \(x=\sum_{a\in A}\lambda_a a\in K\), with \(\lambda_a\ge 0\) and \(\sum_a\lambda_a=1\), then
\[
\langle u,x\rangle=\sum_{a\in A}\lambda_a\langle u,a\rangle\le \sum_{a\in A}\lambda_a m=m.
\]
Hence \(h_K(u)=m\).

Equality holds iff every \(a\) with \(\lambda_a>0\) satisfies \(\langle u,a\rangle=m\), i.e. iff \(x\in \operatorname{conv}(M_A(u))\). So
\[
F_K(u)=\operatorname{conv}(M_A(u)).
\]
\(\square\)

Two consequences will be used repeatedly:

- if one point of \(A\) strictly beats all others for \(\langle u,\cdot\rangle\), then \(F_K(u)\) is that singleton;
- if \(a\in A\) is a maximizer, then \(a\in F_K(u)\).

---

### 1.3. Exposed faces of a Minkowski sum

For nonempty compact convex sets \(A,B\subset \mathbb R^2\),
\[
F_{A+B}(u)=F_A(u)+F_B(u).
\]

#### Proof
First,
\[
h_{A+B}(u)=\max_{a\in A,\ b\in B}\langle u,a+b\rangle
=\max_{a,b}\bigl(\langle u,a\rangle+\langle u,b\rangle\bigr)
=h_A(u)+h_B(u).
\]
If \(a\in F_A(u)\) and \(b\in F_B(u)\), then
\[
\langle u,a+b\rangle=h_A(u)+h_B(u)=h_{A+B}(u),
\]
so
\[
F_A(u)+F_B(u)\subseteq F_{A+B}(u).
\]

Conversely, if \(z\in F_{A+B}(u)\), write \(z=a+b\) with \(a\in A\), \(b\in B\). Then
\[
h_A(u)+h_B(u)=h_{A+B}(u)=\langle u,z\rangle
=\langle u,a\rangle+\langle u,b\rangle\le h_A(u)+h_B(u),
\]
so both inequalities are equalities:
\[
a\in F_A(u),\qquad b\in F_B(u).
\]
Hence
\[
z\in F_A(u)+F_B(u).
\]
Thus
\[
F_{A+B}(u)=F_A(u)+F_B(u).
\]
\(\square\)

---

### 1.4. Exposed faces in the plane

If \(K\subset \mathbb R^2\) is compact, convex, and \(2\)-dimensional, then for every nonzero \(u\),
\[
F_K(u)=K\cap L_u,\qquad L_u:=\{x:\langle u,x\rangle=h_K(u)\}.
\]
Since \(L_u\) is a line, \(F_K(u)\) is a compact convex subset of a line; therefore \(F_K(u)\) is either

- a point, or
- a compact segment.

If \(K\) is a polygon, the point case is a vertex, and the nondegenerate segment case is an edge.

In particular, for our polygon \(P\), every exposed face \(F_P(u)\) is either

- \(\{s_i\}\) for some \(i\), or
- \(E_i=[s_i,s_{i+1}]\) for some \(i\).

---

### 1.5. Compact convex subsets of a line

Let \(a\neq 0\), and let \(G\) be a nonempty compact convex subset of a line parallel to \(a\). If \(p,q\in G\) are respectively the minimum and maximum points of \(\langle a,\cdot\rangle\) on \(G\), then
\[
G=[p,q].
\]

#### Proof
Write the line as \(x_0+\mathbb Ra\). Since \(G\) is compact and convex,
\[
G=\{x_0+\tau a:\tau\in[\alpha,\beta]\}
\]
for some \(\alpha\le \beta\). Then
\[
\langle a,x_0+\tau a\rangle=\langle a,x_0\rangle+\tau\|a\|^2,
\]
which is strictly increasing in \(\tau\). So the minimum occurs at \(\tau=\alpha\), the maximum at \(\tau=\beta\), and therefore \(G=[p,q]\).
\(\square\)

---

### 1.6. Cyclic order of outer normals of a convex polygon

**Lemma.** Let \(P\) be a convex polygon with vertices \(v_1,\dots,v_n\) listed in counterclockwise order (indices mod \(n\)). For each \(i\), let
\[
\alpha_i := v_{i+1}-v_i
\]
be the edge direction and let \(n_i\) be the outer unit normal to the edge \([v_i,v_{i+1}]\), obtained by rotating \(\alpha_i\) by \(-90^\circ\). Then as \(i\) increases from \(1\) to \(n\) cyclically, the directions of \(n_i\) rotate strictly counterclockwise through a full \(360^\circ\). In particular, the map \(i\mapsto n_i\) is a cyclic-order-preserving bijection from \(\{1,\dots,n\}\) to the set of outer normal directions of \(P\).

#### Proof
For each \(i\), let \(\theta_i\) be the exterior angle of \(P\) at vertex \(v_{i+1}\), defined as the angle from direction \(\alpha_i\) to direction \(\alpha_{i+1}\), measured counterclockwise. Two facts are standard:

1. **Strict positivity:** Since \(P\) is a convex polygon, all vertices are extreme points, so no three vertices are collinear. At each vertex \(v_{i+1}\) the boundary of \(P\) turns strictly to the left, giving \(\theta_i>0\).

2. **Total turning:** The total turning angle of a closed convex polygon traversed counterclockwise equals exactly \(2\pi\):
\[
\sum_{i=1}^{n}\theta_i=2\pi.
\]

Since \(n_i\) is \(\alpha_i\) rotated by the fixed amount \(-90^\circ\), the angle from \(n_i\) to \(n_{i+1}\) (measured counterclockwise) equals the angle from \(\alpha_i\) to \(\alpha_{i+1}\), which is \(\theta_i\). Therefore:

- At each step \(i\to i+1\), the normal \(n_{i+1}\) is obtained from \(n_i\) by a counterclockwise rotation of \(\theta_i>0\).
- After all \(n\) steps, the total counterclockwise rotation is \(\sum_{i=1}^n\theta_i=2\pi\), returning to the starting direction.

Hence the sequence \(n_1,n_2,\dots,n_n\) traces out distinct directions that appear in strictly increasing counterclockwise order, cycling once around the full circle. This means:

- All \(n_i\) are distinct (the map \(i\mapsto n_i\) is injective).
- The cyclic order of the indices \(1,\dots,n\) is preserved by the map \(i\mapsto n_i\) (when directions are ordered counterclockwise around the circle).

This is precisely the statement that \(i\mapsto n_i\) is a cyclic-order-preserving bijection.
\(\square\)

---

## 2. Proof of statement (2)

Fix \(i\), and let \(u_i\) be an outer normal to the edge
\[
E_i=[s_i,s_{i+1}].
\]
Then
\[
F_P(u_i)=E_i,
\]
so
\[
\langle u_i,s_i\rangle=\langle u_i,s_{i+1}\rangle=h_P(u_i),
\]
and for every \(j\neq i,i+1\),
\[
\langle u_i,s_j\rangle<h_P(u_i).
\]

Let
\[
G_i:=F_Q(u_i).
\]

We will prove
\[
G_i=[t_i,t_{i+1}].
\]

### Step 1: small perturbations of \(u_i\) enter the adjacent cones

For \(j\neq i,i+1\), set
\[
\delta_j:=\langle u_i,s_i-s_j\rangle=\langle u_i,s_{i+1}-s_j\rangle>0,
\qquad
\delta:=\min_{j\neq i,i+1}\delta_j>0.
\]
Also set
\[
M:=\max_{j\neq i,i+1}\max\Bigl(
|\langle a_i,s_i-s_j\rangle|,\ |\langle a_i,s_{i+1}-s_j\rangle|
\Bigr).
\]
Choose
\[
0<\varepsilon<\frac{\delta}{M+1}.
\]

Because \(u_i\perp a_i\), for \(j\neq i,i+1\) we have
\[
\langle u_i-\varepsilon a_i,s_i-s_j\rangle
=\delta_j-\varepsilon\langle a_i,s_i-s_j\rangle
\ge \delta-\varepsilon M>0,
\]
and
\[
\langle u_i-\varepsilon a_i,s_i-s_{i+1}\rangle
=-\varepsilon\langle a_i,s_i-s_{i+1}\rangle
=\varepsilon\|a_i\|^2>0.
\]
So \(s_i\) strictly beats every other vertex of \(P\) for the functional \(\langle u_i-\varepsilon a_i,\cdot\rangle\). By §1.2,
\[
F_P(u_i-\varepsilon a_i)=\{s_i\},
\]
hence
\[
u_i-\varepsilon a_i\in C_i.
\]

Similarly,
\[
\langle u_i+\varepsilon a_i,s_{i+1}-s_j\rangle\ge \delta-\varepsilon M>0
\quad (j\neq i,i+1),
\]
and
\[
\langle u_i+\varepsilon a_i,s_{i+1}-s_i\rangle
=\varepsilon\|a_i\|^2>0.
\]
Therefore
\[
F_P(u_i+\varepsilon a_i)=\{s_{i+1}\},
\]
so
\[
u_i+\varepsilon a_i\in C_{i+1}.
\]

---

### Step 2: \(t_i,t_{i+1}\in G_i\)

Take a sequence \(\varepsilon_m\downarrow 0\) such that
\[
u_i-\varepsilon_m a_i\in C_i,\qquad u_i+\varepsilon_m a_i\in C_{i+1}
\]
for every \(m\).

Since \(t_i\) is the unique maximizer of every direction in \(C_i\), for every \(t\in T\),
\[
\langle u_i-\varepsilon_m a_i,\ t_i-t\rangle\ge 0.
\]
Letting \(m\to\infty\) gives
\[
\langle u_i,\ t_i-t\rangle\ge 0\qquad(\forall t\in T).
\]
Thus \(t_i\) maximizes \(\langle u_i,\cdot\rangle\) on \(T\), so by §1.2,
\[
t_i\in F_Q(u_i)=G_i.
\]

The same argument with \(C_{i+1}\) gives
\[
t_{i+1}\in G_i.
\]

---

### Step 3: \(G_i\) is exactly \([t_i,t_{i+1}]\)

Because
\[
G_i=F_Q(u_i)\subset \{q:\langle u_i,q\rangle=h_Q(u_i)\},
\]
the set \(G_i\) lies in a line orthogonal to \(u_i\), hence parallel to \(a_i\).

Now fix \(\varepsilon>0\) as in Step 1. For \(q\in G_i\), we have \(\langle u_i,q\rangle=h_Q(u_i)\), so
\[
\langle u_i-\varepsilon a_i,q\rangle=h_Q(u_i)-\varepsilon\langle a_i,q\rangle.
\]
Therefore, **on \(G_i\)**, maximizing \(\langle u_i-\varepsilon a_i,\cdot\rangle\) is equivalent to minimizing \(\langle a_i,\cdot\rangle\).

But \(u_i-\varepsilon a_i\in C_i\), and \(t_i\) is the unique maximizer of that direction on \(T\). By §1.2,
\[
F_Q(u_i-\varepsilon a_i)=\{t_i\}.
\]
This means \(t_i\) is the unique maximizer of \(\langle u_i-\varepsilon a_i,\cdot\rangle\) on all of \(Q\). Since \(G_i\subseteq Q\) and \(t_i\in G_i\) (established in Step 2), \(t_i\) is therefore the unique maximizer of \(\langle u_i-\varepsilon a_i,\cdot\rangle\) on \(G_i\). Combined with the equivalence above (on \(G_i\), maximizing \(\langle u_i-\varepsilon a_i,\cdot\rangle\) is equivalent to minimizing \(\langle a_i,\cdot\rangle\)), it follows that \(t_i\) is the unique minimizer of \(\langle a_i,\cdot\rangle\) on \(G_i\).

Likewise, because \(u_i+\varepsilon a_i\in C_{i+1}\),
\[
F_Q(u_i+\varepsilon a_i)=\{t_{i+1}\},
\]
and on \(G_i\),
\[
\langle u_i+\varepsilon a_i,q\rangle=h_Q(u_i)+\varepsilon\langle a_i,q\rangle,
\]
so \(t_{i+1}\) is the unique maximizer of \(\langle a_i,\cdot\rangle\) on \(G_i\).

Now \(G_i\) is a compact convex subset of a line parallel to \(a_i\), and \(t_i,t_{i+1}\) are respectively the minimum and maximum points of \(\langle a_i,\cdot\rangle\) on \(G_i\). By §1.5,
\[
G_i=[t_i,t_{i+1}].
\]

Since multiplying \(u_i\) by a positive scalar does not change the exposed face, this is true for **any** outer normal to \(E_i\).

---

### Step 4: direction of \(t_{i+1}-t_i\)

Because \(G_i=[t_i,t_{i+1}]\) lies in a line parallel to \(a_i\), there exists \(\lambda_i\in\mathbb R\) such that
\[
t_{i+1}-t_i=\lambda_i a_i=\lambda_i(s_{i+1}-s_i).
\]
Also \(t_i\) is the minimum and \(t_{i+1}\) the maximum of \(\langle a_i,\cdot\rangle\) on \(G_i\), so
\[
\langle a_i,t_{i+1}-t_i\rangle\ge 0.
\]
Substituting \(t_{i+1}-t_i=\lambda_i a_i\) gives
\[
\lambda_i\|a_i\|^2\ge 0.
\]
Since \(a_i\neq 0\), we obtain
\[
\lambda_i\ge 0.
\]

This proves statement (2).

---

## 3. Proof of statement (1)

Recall
\[
x_i=s_i+t_i.
\]

### Step 1: each \(x_i\) is a vertex of \(R\)

If \(u\in C_i\), then by definition of \(C_i\),
\[
F_P(u)=\{s_i\}.
\]
By the Subproblem 1–2 input, \(t_i\) is the unique maximizer of that same direction on \(T\); by §1.2,
\[
F_Q(u)=\{t_i\}.
\]
Hence by §1.3,
\[
F_R(u)=F_P(u)+F_Q(u)=\{s_i+t_i\}=\{x_i\}.
\]
This is a singleton exposed face of the polygon \(R\), so \(x_i\) is a vertex of \(R\).

---

### Step 2: for each \(i\), \([x_i,x_{i+1}]\) is an edge of \(R\)

From statement (2),
\[
F_Q(u_i)=[t_i,t_{i+1}],
\qquad
t_{i+1}-t_i=\lambda_i a_i
\quad\text{with }\lambda_i\ge 0.
\]
Therefore
\[
E_i=s_i+[0,1]a_i,
\qquad
[t_i,t_{i+1}]=t_i+[0,\lambda_i]a_i.
\]
Using §1.3,
\[
F_R(u_i)=F_P(u_i)+F_Q(u_i)=E_i+[t_i,t_{i+1}]
= x_i+[0,1+\lambda_i]a_i.
\]
But
\[
x_{i+1}
=s_{i+1}+t_{i+1}
=(s_i+a_i)+(t_i+\lambda_i a_i)
=x_i+(1+\lambda_i)a_i.
\]
So
\[
F_R(u_i)=[x_i,x_{i+1}].
\]
Since \(1+\lambda_i>0\) and \(a_i\neq 0\), this segment is nondegenerate. Hence it is an edge of \(R\).

Let me denote it by
\[
e_i:=[x_i,x_{i+1}]=F_R(u_i).
\]

---

### Step 3: these are all the edges of \(R\)

Let \(e\) be any edge of \(R\), and let \(u\) be an outer normal to \(e\). Then
\[
e=F_R(u).
\]

By §1.4, \(F_P(u)\) is either a point or a segment.

- If \(F_P(u)\) is a point, then it must be \(\{s_i\}\) for some \(i\). Hence \(u\in C_i\), so
  \[
  F_Q(u)=\{t_i\}
  \]
  and therefore
  \[
  F_R(u)=\{x_i\},
  \]
  a point. But \(e\) is an edge, so this is impossible.

- Therefore \(F_P(u)\) is a nondegenerate segment. Since the edges of \(P\) are exactly the \(E_i\), we must have
  \[
  F_P(u)=E_i
  \]
  for some \(i\). Then statement (2) gives
  \[
  F_Q(u)=[t_i,t_{i+1}],
  \]
  and so
  \[
  e=F_R(u)=E_i+[t_i,t_{i+1}]=[x_i,x_{i+1}]=e_i.
  \]

Thus every edge of \(R\) is one of
\[
e_1,e_2,\dots,e_n.
\]

---

### Step 4: distinctness and cyclic order

We now justify carefully that the vertices are exactly \(x_1,\dots,x_n\) in cyclic order.

First, the edges \(e_i\) are pairwise distinct. Indeed, each nondegenerate edge has a unique outer-normal ray. The outer-normal ray of \(e_i\) is the same as that of \(u_i\). If \(e_i=e_j\), then \(u_i\) and \(u_j\) would determine the same outer-normal ray. But then that ray would expose both \(E_i\) and \(E_j\) on \(P\), impossible unless \(E_i=E_j\), i.e. \(i=j\).

So \(R\) has exactly \(n\) edges:
\[
e_1,e_2,\dots,e_n.
\]

Now, the boundary of a convex polygon is a simple closed polygonal chain. Hence a polygon with \(n\) edges has exactly \(n\) vertices, namely the intersections of consecutive edges in boundary order.

By Lemma §1.6 applied to the convex polygon \(P\) (whose vertices \(s_1,\dots,s_n\) are in counterclockwise order), the outer normal directions of the edges \(E_1,\dots,E_n\) appear in strictly increasing counterclockwise order around the circle: that is, the cyclic order of the outer-normal rays \(u_1,u_2,\dots,u_n\) matches the cyclic order of the edge indices \(1,2,\dots,n\). Since \(e_i\) has the same outer-normal ray as \(E_i\), it follows that the edges \(e_1,e_2,\dots,e_n\) of \(R\) have outer normals in the same strict counterclockwise cyclic order.

Now we need to conclude that the edges \(e_1,\dots,e_n\) of \(R\) occur on the boundary of \(R\) in that cyclic order. Apply Lemma §1.6 to the convex polygon \(R\) itself: it gives a cyclic-order-preserving bijection \(\varphi\colon\{1,\dots,n\}\to\{1,\dots,n\}\) such that the boundary order of the edges of \(R\) is \(e_{\varphi(1)},e_{\varphi(2)},\dots,e_{\varphi(n)}\) and the corresponding outer normals \(u_{\varphi(1)},u_{\varphi(2)},\dots,u_{\varphi(n)}\) appear in strictly increasing counterclockwise order. We have already shown that \(u_1,u_2,\dots,u_n\) also appear in strictly increasing counterclockwise order. Since the \(n\) directions \(u_1,\dots,u_n\) are pairwise distinct (established above, as \(e_i\) are distinct edges with distinct outer normals), a set of \(n\) distinct directions on the circle admits a unique counterclockwise cyclic order. Therefore the two sequences \((u_{\varphi(1)},\dots,u_{\varphi(n)})\) and \((u_1,\dots,u_n)\) represent the same cyclic order, which means \(\varphi\) is a cyclic rotation: \(\varphi(k)=k+c\pmod n\) for some constant \(c\). In particular, the cyclic order \((e_{\varphi(1)},\dots,e_{\varphi(n)})=(e_{1+c},\dots,e_{n+c})\) is the same cyclic order as \((e_1,\dots,e_n)\). Hence the edges
\[
e_1,e_2,\dots,e_n
\]
occur on the boundary of \(R\) in that same cyclic order.

Therefore the vertices of \(R\) are
\[
e_n\cap e_1,\ e_1\cap e_2,\ \dots,\ e_{n-1}\cap e_n.
\]
But by definition,
\[
x_i\in e_{i-1}\cap e_i
\]
for every \(i\). Since consecutive edges of a polygon meet in exactly one vertex, we get
\[
e_{i-1}\cap e_i=\{x_i\}.
\]
Thus the boundary vertices of \(R\) are exactly
\[
x_1,x_2,\dots,x_n
\]
in cyclic order. In particular, they are pairwise distinct.

This proves statement (1).

---

## Conclusion

For each \(i\),

1. the point
   \[
   x_i=s_i+t_i
   \]
   is a vertex of \(P+Q\), and the vertices of \(P+Q\) are exactly
   \[
   x_1,x_2,\dots,x_n
   \]
   in cyclic order;

2. for any outer normal \(u\) to the edge \([s_i,s_{i+1}]\) of \(P\),
   \[
   F_Q(u)=[t_i,t_{i+1}],
   \]
   and therefore
   \[
   t_{i+1}-t_i=\lambda_i(s_{i+1}-s_i)
   \qquad\text{for some }\lambda_i\ge 0.
   \]

So both requested statements are proved.
